\section*{Source Notes / この号の読み方}

本誌はニュース記事の引用集ではなく、\textbf{event chronology、technical evidence、community observationを分離したsurvey}として作成している。読者が各claimの強さを読み分けられるよう、本文では次の4種類を区別する。

\begin{tcolorbox}[enhanced,breakable,colback=SurveySoft,colframe=SurveyRule,boxrule=0.4pt,arc=1mm,left=3mm,right=3mm,top=2mm,bottom=2mm]
\textbf{Primary / Technical source}\par
Official announcement、documentation、model card、paper、GitHub release等。存在・chronology・仕様の核は原則としてここで確認する。Vendor / author自身のbenchmark値は、primary sourceであっても独立検証済み性能とは区別する。\par\smallskip
\textbf{Claim Boundary}\par
Vendor claim、author-reported metric、simulation result、benchmark条件、未検証事項など、主張の適用範囲を本文から切り分けて表示する。\par\smallskip
\textbf{Community Observation}\par
GrokをX trend sensorとして使い、具体的なpost URLをReaction Passで収集した。「その反応がX上で観測された」ことのEvidenceであり、post内のtechnical claimの正しさを保証しない。\par\smallskip
\textbf{Late Breaking}\par
Friday 18:00 America/New\_YorkのEditorial Cutoff後に重要度が上がった事項。Main-window eventへ逆流させず、短報として分離する。
\end{tcolorbox}

今号の候補母集団は35件。Candidate Inventoryからprimary verification、paper full review、X Reaction Evidence、cross-candidate comparisonを経て掲載対象を選定した。誤ったchronologyや一次情報で確認できなかった候補も削除せず、Repository内にprovenanceとして保持している。

各citationは、その直前の具体的claimを支えるsourceへ接続する。SOTA、benchmark superiority、resource requirement等は、setupやattributionを伴わず一般化しない。詳細なscreening recordとEvidence ReviewはRepository \texttt{eariver/japanese-generative-ai-survey} の \texttt{sources/2026-W32/} 以下に保存している\cite{repo2026w32}。
